\documentclass[10pt,letterpaper,twocolumn]{article}

\usepackage[
  letterpaper,
  top=0.7in,
  bottom=0.7in,
  left=0.65in,
  right=0.65in,
  columnsep=0.2in
]{geometry}
\usepackage[T1]{fontenc}
\usepackage{mathptmx}
\usepackage{amsmath}
\usepackage{booktabs}
\usepackage{array}
\usepackage{cite}
\usepackage{url}
\usepackage{float}

\setlength{\parindent}{0.14in}
\setlength{\parskip}{0pt}
\pagestyle{plain}
\renewcommand{\thesection}{\Roman{section}}
\renewcommand{\thesubsection}{\Alph{subsection}.}
\renewcommand{\thetable}{\Roman{table}}
\renewcommand{\refname}{REFERENCES}

\makeatletter
\renewcommand\section{\@startsection{section}{1}{\z@}%
  {1.3ex plus 1ex minus .2ex}%
  {0.8ex plus .2ex}%
  {\centering\normalfont\small\scshape}}
\renewcommand\subsection{\@startsection{subsection}{2}{\z@}%
  {1.0ex plus .8ex minus .2ex}%
  {0.6ex plus .2ex}%
  {\normalfont\itshape}}
\renewcommand{\fnum@table}{\textsc{Table \thetable}}
\makeatother

\begin{document}

\twocolumn[
\begin{@twocolumnfalse}
\begin{center}
{\LARGE Lightweight RGB Hand Pose Tracking for Edge HCI\par}
\vspace{0.6em}
{\large Brandon MacGillivray\par}
\vspace{0.2em}
{\normalsize School of Electronics, Electrical Engineering and Computer Science\par}
{\normalsize Queen's University Belfast, Belfast, UK\par}
{\normalsize \texttt{bmacgillivray01@qub.ac.uk}\par}
\end{center}

\vspace{0.6em}
\noindent\textbf{\textit{Abstract}}---This stripped-down version keeps the paper structure and results while reducing the prose to core points only. The work evaluates a DRHand-inspired RGB hand pose pipeline through RHD ablations, transfer-learning experiments, a MediaPipe baseline comparison, and Jetson Orin Nano benchmarking. The main findings are that a reduced 10-keypoint model preserves most shared-subset accuracy while reducing latency, fusion is generally stronger than either branch alone, target-domain fine-tuning helps transfer substantially, and the MediaPipe baseline is weaker on RHD and sensitive to handedness filtering.

\vspace{0.6em}
\noindent\textbf{\textit{Keywords}}---hand pose estimation, human-computer interaction, edge devices, lightweight deep learning, stereo vision
\vspace{1.0em}
\end{@twocolumnfalse}
]

\section{INTRODUCTION}

Hand pose tracking is relevant to HCI because it supports touchless interaction in settings where direct input is inconvenient or undesirable \cite{vuletic,wachs,li_vr}. For edge deployment, the central problem is the trade-off between geometric accuracy and runtime cost on constrained hardware \cite{mediapipe,dualreg,srhandnet,fasthand,attentionlite}. This paper studies that trade-off for a lightweight DRHand-inspired RGB pipeline on an embedded Jetson target.

The main points are:
\begin{itemize}
\item A 21-keypoint baseline is compared with a reduced 10-keypoint variant.
\item Fusion is compared with heatmap-only and coordinate-only prediction.
\item Transfer is evaluated between RHD and a COCO-style hand dataset.
\item A MediaPipe baseline is included for additional context.
\item Deployment cost is measured on Jetson Orin Nano.
\end{itemize}

\section{LITERATURE REVIEW}

The relevant literature spans touchless HCI, lightweight pose estimation, and embedded deployment. Gesture-based HCI work motivates hand tracking as a practical input modality \cite{vuletic,wachs,li_vr}. MediaPipe Hands and related on-device systems show that efficient landmark prediction is feasible on commodity hardware \cite{mediapipe,blazepose}. DRHand is the closest methodological reference because it combines heatmap and coordinate regression in a lightweight dual-branch architecture \cite{dualreg}. Stereo and embedded pose studies motivate the wider project direction beyond the present RGB baseline \cite{stereojetson}.

\section{TECHNICAL APPROACH}

This section outlines the scope, model, experimental design, and evaluation criteria used in the study.

\subsection{Project Scope and Motivation}

The central design question is whether all 21 hand keypoints are necessary for edge HCI tasks. A reduced 10-keypoint layout based on finger bases and tips may retain useful interaction information while lowering output complexity and runtime. The study therefore compares 21-joint and 10-joint variants and extends the analysis to transfer and baseline comparison.

\subsection{DRHand-Inspired Model}

The implemented model uses a shared lightweight backbone with two output branches:
\begin{itemize}
\item a heatmap branch for spatial localisation
\item a coordinate branch for direct geometric regression
\end{itemize}

The branches are trained with separate losses and combined during inference by a simple reliability rule. When the branches agree closely, the heatmap prediction is used; otherwise the coordinate prediction is preferred. This design follows the central idea of DRHand while adapting it to the current repository \cite{dualreg,wingloss}.

\subsection{Relation to the Original DRHand Paper}

The implementation should be read as an adaptation, not a strict reproduction. The main differences are:
\begin{itemize}
\item different dataset setup
\item PyTorch rather than TensorFlow
\item configurable schedules with early stopping
\item gradient accumulation in stage 2
\item added transfer-learning and 10-keypoint functionality
\end{itemize}

These differences mean that direct numeric comparison with the original DRHand paper is not the main goal.

\subsection{Experimental Setup}

The result set contains four experimental components:
\begin{itemize}
\item five-seed RHD ablation experiments
\item transfer-learning runs between RHD and COCO-style hand data
\item MediaPipe baseline evaluation on RHD
\item Jetson Orin Nano benchmarking
\end{itemize}

Training follows a two-stage procedure. Stage 1 trains the backbone and heatmap branch. Stage 2 trains the coordinate branch jointly with the rest of the model. Evaluation is performed in fusion, heatmap-only, and coordinate-only modes. Shared-10 evaluation is used to compare the 21-joint and 10-joint models fairly on the same subset.

\subsection{Evaluation Metrics}

The main accuracy metrics are normalised SSE, normalised EPE, and PCK@0.2. Additional fusion diagnostics record branch-selection behaviour and oracle-match rate. Benchmarking records forward time, end-to-end time, images per second, and detailed pipeline timing. These measures support both accuracy analysis and deployment analysis.

\section{RESULTS AND ANALYSIS}

This section summarises the main findings from the ablation, transfer, baseline, and benchmarking results.

\subsection{Main Accuracy Results}

Table~\ref{tab:main_accuracy} is the core accuracy comparison. The key points are:
\begin{itemize}
\item \texttt{B0} fusion on native-21 gives SSE 0.3098, EPE 0.0779, and PCK 0.9112.
\item \texttt{B0} shared-10 gives SSE 0.1511, EPE 0.0764, and PCK 0.9101.
\item \texttt{B1} native-10 gives SSE 0.1574, EPE 0.0706, and PCK 0.9099.
\end{itemize}

The fair comparison is \texttt{B0} shared-10 against \texttt{B1} native-10. On that comparison, the two models are very close: \texttt{B0} is slightly better on SSE, \texttt{B1} is slightly better on EPE, and PCK is effectively tied.

\subsection{Branch Ablation}

Table~\ref{tab:branch_ablation} shows that fusion is the strongest overall operating mode. For \texttt{B0}, heatmap-only gives the highest PCK but clearly worse EPE than fusion. For \texttt{B1}, heatmap-only degrades more noticeably across the metrics. The result supports the claim that the heatmap and coordinate branches are complementary rather than redundant.

\subsection{Transfer Learning}

Table~\ref{tab:transfer_learning} shows three main patterns:
\begin{itemize}
\item zero-shot cross-dataset generalisation is weak
\item target-domain fine-tuning helps strongly in both directions
\item target-domain scratch remains strongest on its home dataset
\end{itemize}

The transfer is asymmetric. \texttt{R\_TO\_C} approaches COCO scratch performance closely, whereas \texttt{C\_TO\_R} remains more clearly below RHD scratch. In the present setup, adaptation from RHD to COCO is therefore easier than the reverse.

\subsection{MediaPipe Baseline}

Table~\ref{tab:mediapipe_baseline} shows that the MediaPipe baseline is substantially weaker than the trained DRHand variants on RHD. The learned models are much stronger on PCK and SSE in both native and shared-10 comparisons. The MediaPipe results also show a strong dependence on handedness filtering: ignoring handedness raises detection success from 0.4879 to 0.8933 and improves PCK, although EPE worsens slightly.

\subsection{Runtime Analysis}

Table~\ref{tab:latency} shows that the 10-joint model is faster than the 21-joint baseline across all three prediction modes. In fusion mode, \texttt{B1} reduces end-to-end time from 36.81 ms/image to 33.31 ms/image and increases throughput from 27.20 to 30.03 images/s. The speedup is useful but modest, which is consistent with the shared backbone remaining the main computational cost.

Table~\ref{tab:latency_breakdown} shows that forward prediction dominates the runtime. However, read, preprocessing, host-to-device transfer, and JSON writing still contribute enough overhead to matter on an embedded interactive system.

\subsection{Fusion Diagnostics}

Table~\ref{tab:fusion_diagnostics} shows that the current hand-crafted fusion heuristic is not close to oracle behaviour. The native oracle-match rates are 0.481 for \texttt{B0} and 0.529 for \texttt{B1}. The available shared-10 diagnostic row for \texttt{B0} remains similarly low at 0.474. Fusion improves aggregate accuracy, but the branch-selection rule is still an obvious target for refinement.

\subsection{Hyperparameter Sensitivity}

Table~\ref{tab:hyper_ablation} shows that no single setting dominates every metric. The main extremes are:
\begin{itemize}
\item \texttt{W2}: best SSE
\item \texttt{E1}: best EPE
\item \texttt{S2}: best PCK
\end{itemize}

Latency differences across these fused 21-joint settings are small. The main implication is that configuration choice depends on which metric matters most.

\section{CONCLUSION}

The main conclusions are:
\begin{itemize}
\item the reduced 10-joint model preserves most useful shared-subset accuracy while lowering latency
\item fusion is generally stronger than either branch alone
\item transfer fine-tuning is effective, but zero-shot cross-dataset generalisation is weak
\item the MediaPipe baseline is much weaker on RHD and is sensitive to handedness filtering
\item the current fusion heuristic still leaves clear room for improvement
\end{itemize}

The pipeline already operates close to real time on embedded hardware. The next improvements should focus on better fusion, broader evaluation, and richer sensing comparisons.

\clearpage
\onecolumn
\begin{thebibliography}{00}

\bibitem{vuletic}
T. Vuletic, A. Duffy, L. Hay, C. McTeague, and G. Campbell, ``Systematic literature review of hand gestures used in human-computer interaction interfaces,'' \emph{Int. J. Human-Computer Studies}, vol. 129, pp. 47--65, 2019.

\bibitem{wachs}
J. P. Wachs, M. Kolsch, H. Stern, and Y. Edan, ``Vision-based hand-gesture applications,'' \emph{Communications of the ACM}, vol. 54, no. 2, pp. 60--71, 2011.

\bibitem{li_vr}
Y. Li, J. Huang, F. Tian, H.-A. Wang, and G.-Z. Dai, ``Gesture interaction in virtual reality,'' \emph{Virtual Reality \& Intelligent Hardware}, vol. 1, no. 1, pp. 84--112, 2019.

\bibitem{dualreg}
D. Wei, S. An, X. Zhang, J. Tian, K. A. Tsintotas, A. Gasteratos, and H. Zhu, ``Dual regression for efficient hand pose estimation,'' in \emph{Proc. IEEE Int. Conf. Robotics and Automation}, 2022.

\bibitem{mediapipe}
F. Zhang, V. Bazarevsky, A. Vakunov, A. Tkachenka, G. Sung, C.-L. Chang, and M. Grundmann, ``MediaPipe hands: On-device real-time hand tracking,'' \emph{arXiv preprint arXiv:2006.10214}, 2020.

\bibitem{srhandnet}
Y. Wang, B. Zhang, and C. Peng, ``SRHandNet: Real-time 2D hand pose estimation with simultaneous region localization,'' \emph{IEEE Trans. Image Processing}, vol. 29, pp. 2977--2986, 2019.

\bibitem{wingloss}
Z.-H. Feng, J. Kittler, M. Awais, P. Huber, and X.-J. Wu, ``Wing loss for robust facial landmark localisation with convolutional neural networks,'' in \emph{Proc. IEEE Conf. Computer Vision and Pattern Recognition}, 2018, pp. 2235--2245.

\bibitem{fasthand}
S. An, X. Zhang, D. Wei, H. Zhu, J. Yang, and K. A. Tsintotas, ``FastHand: Fast monocular hand pose estimation on embedded systems,'' \emph{Journal of Systems Architecture}, vol. 122, p. 102361, 2022.

\bibitem{attentionlite}
N. Santavas, I. Kansizoglou, L. Bampis, E. Karakasis, and A. Gasteratos, ``Attention! A lightweight 2D hand pose estimation approach,'' \emph{IEEE Sensors Journal}, vol. 21, no. 10, pp. 11488--11501, 2021.

\bibitem{stereojetson}
E. Williams-Linera and J. M. Ramirez-Cortes, ``Stereo vision system based on the NVIDIA Jetson Nano for real-time evaluation of yoga poses,'' in \emph{Proc. IEEE}, 2024, pp. 1--7.

\bibitem{blazepose}
V. Bazarevsky, I. Grishchenko, K. Raveendran, T. Zhu, F. Zhang, and M. Grundmann, ``BlazePose: On-device real-time body pose tracking,'' \emph{arXiv preprint arXiv:2006.10204}, 2020.

\end{thebibliography}

\clearpage
\onecolumn
\section*{TABLES}

\begin{table}[H]
\centering
\footnotesize
\caption{Main quantitative comparison between the 21-keypoint baseline and the reduced 10-keypoint model on RHD. Lower SSE and EPE are better; higher PCK@0.2 is better. Values are mean $\pm$ std over five seeds.}
\label{tab:main_accuracy}
\begin{tabular}{llcccc}
\toprule
Model & Eval set & Train joints & SSE $\downarrow$ & EPE $\downarrow$ & PCK@0.2 $\uparrow$ \\
\midrule
B0 baseline & native-21 & 21 & 0.3098 $\pm$ 0.0236 & 0.0779 $\pm$ 0.0054 & 0.9112 $\pm$ 0.0071 \\
B0 baseline & shared-10 & 21 & 0.1511 $\pm$ 0.0120 & 0.0764 $\pm$ 0.0039 & 0.9101 $\pm$ 0.0071 \\
B1 reduced & native-10 & 10 & 0.1574 $\pm$ 0.0226 & 0.0706 $\pm$ 0.0054 & 0.9099 $\pm$ 0.0046 \\
\bottomrule
\end{tabular}
\end{table}

\begin{table}[H]
\centering
\footnotesize
\caption{Branch ablation for the dual-regression design. Fusion is compared against the heatmap-only and coordinate-only outputs using the same trained checkpoints. Values are mean $\pm$ std over five seeds.}
\label{tab:branch_ablation}
\begin{tabular}{llcccc}
\toprule
Model & Eval set & Mode & SSE $\downarrow$ & EPE $\downarrow$ & PCK@0.2 $\uparrow$ \\
\midrule
B0 & native-21 & fusion & 0.3098 $\pm$ 0.0236 & 0.0779 $\pm$ 0.0054 & 0.9112 $\pm$ 0.0071 \\
B0 & native-21 & heatmap & 0.3139 $\pm$ 0.0684 & 0.0989 $\pm$ 0.0335 & 0.9374 $\pm$ 0.0066 \\
B0 & native-21 & coord & 0.3134 $\pm$ 0.0229 & 0.0788 $\pm$ 0.0062 & 0.9113 $\pm$ 0.0071 \\
B1 & native-10 & fusion & 0.1574 $\pm$ 0.0226 & 0.0706 $\pm$ 0.0054 & 0.9099 $\pm$ 0.0046 \\
B1 & native-10 & heatmap & 0.2213 $\pm$ 0.0880 & 0.0991 $\pm$ 0.0288 & 0.9022 $\pm$ 0.0393 \\
B1 & native-10 & coord & 0.1599 $\pm$ 0.0218 & 0.0715 $\pm$ 0.0070 & 0.9102 $\pm$ 0.0048 \\
\bottomrule
\end{tabular}
\end{table}

\begin{table}[H]
\centering
\footnotesize
\caption{Transfer-learning results for the four training sequences evaluated on both RHD and the COCO-style hand dataset. Values are mean $\pm$ std over five seeds, and fusion-mode evaluation is shown because it is the main operating mode in the study.}
\label{tab:transfer_learning}
\begin{tabular}{llcccccc}
\toprule
Training sequence & Exp. & RHD SSE & RHD EPE & RHD PCK & COCO SSE & COCO EPE & COCO PCK \\
\midrule
RHD scratch & \texttt{R\_ONLY} & 0.3091 $\pm$ 0.0269 & 0.0764 $\pm$ 0.0030 & 0.9143 $\pm$ 0.0097 & 3.7488 $\pm$ 0.7985 & 0.3365 $\pm$ 0.0075 & 0.3719 $\pm$ 0.0434 \\
COCO scratch & \texttt{C\_ONLY} & 1.4194 $\pm$ 0.2619 & 0.1609 $\pm$ 0.0207 & 0.3397 $\pm$ 0.0718 & 0.0764 $\pm$ 0.0016 & 0.0612 $\pm$ 0.0049 & 0.9834 $\pm$ 0.0007 \\
COCO $\rightarrow$ RHD & \texttt{C\_TO\_R} & 0.4504 $\pm$ 0.0162 & 0.1020 $\pm$ 0.0039 & 0.8190 $\pm$ 0.0114 & 1.6426 $\pm$ 0.2293 & 0.2678 $\pm$ 0.0114 & 0.5320 $\pm$ 0.0164 \\
RHD $\rightarrow$ COCO & \texttt{R\_TO\_C} & 0.5699 $\pm$ 0.0592 & 0.1276 $\pm$ 0.0149 & 0.7690 $\pm$ 0.0221 & 0.0980 $\pm$ 0.0064 & 0.0840 $\pm$ 0.0018 & 0.9837 $\pm$ 0.0007 \\
\bottomrule
\end{tabular}
\end{table}

\begin{table}[H]
\centering
\footnotesize
\caption{Comparison between the MediaPipe Hand Landmarker baseline and the learned DRHand variants on RHD. The MediaPipe rows are single evaluation runs, whereas the DRHand rows are mean $\pm$ std over five seeds. Detection success is shown only for the MediaPipe runs because the learned-model evaluations do not use the same detection stage.}
\label{tab:mediapipe_baseline}
\begin{tabular}{llcccc}
\toprule
Method & Eval set & SSE & EPE & PCK@0.2 & Detect. succ. \\
\midrule
MediaPipe & native-21 & 2.3022 & 0.0830 & 0.6452 & 0.4879 \\
MediaPipe (ignore handedness) & native-21 & 1.4002 & 0.0904 & 0.6726 & 0.8933 \\
MediaPipe & shared-10 & 1.1202 & 0.1209 & 0.6407 & 0.4879 \\
\texttt{B0} fusion & native-21 & 0.3098 $\pm$ 0.0236 & 0.0779 $\pm$ 0.0054 & 0.9112 $\pm$ 0.0071 & n/a \\
\texttt{B0} fusion & shared-10 & 0.1511 $\pm$ 0.0120 & 0.0764 $\pm$ 0.0039 & 0.9101 $\pm$ 0.0071 & n/a \\
\texttt{B1} fusion & native-10 & 0.1574 $\pm$ 0.0226 & 0.0706 $\pm$ 0.0054 & 0.9099 $\pm$ 0.0046 & n/a \\
\bottomrule
\end{tabular}
\end{table}

\begin{table}[H]
\centering
\footnotesize
\caption{Edge-device latency summary for the main models. Forward time measures the prediction stage only, whereas end-to-end time includes image read, preprocessing, host-to-device transfer, prediction, and JSON write. Values are mean $\pm$ std over five seeds.}
\label{tab:latency}
\begin{tabular}{llccc}
\toprule
Model & Mode & Forward ms/img $\downarrow$ & End-to-end ms/img $\downarrow$ & Images/s $\uparrow$ \\
\midrule
B0 & fusion & 27.43 $\pm$ 0.45 & 36.81 $\pm$ 1.33 & 27.20 $\pm$ 0.97 \\
B0 & heatmap & 24.28 $\pm$ 0.35 & 33.19 $\pm$ 0.56 & 30.13 $\pm$ 0.50 \\
B0 & coord & 24.08 $\pm$ 0.44 & 33.17 $\pm$ 0.86 & 30.16 $\pm$ 0.77 \\
B1 & fusion & 24.59 $\pm$ 0.61 & 33.31 $\pm$ 0.71 & 30.03 $\pm$ 0.64 \\
B1 & heatmap & 24.00 $\pm$ 0.34 & 32.64 $\pm$ 0.46 & 30.64 $\pm$ 0.43 \\
B1 & coord & 23.91 $\pm$ 0.29 & 32.50 $\pm$ 0.34 & 30.77 $\pm$ 0.32 \\
\bottomrule
\end{tabular}
\end{table}

\begin{table}[H]
\centering
\footnotesize
\caption{Detailed latency breakdown for the main fused models. These metrics are available from the benchmark pipeline and help localize where runtime is spent. Values are mean $\pm$ std over five seeds.}
\label{tab:latency_breakdown}
\begin{tabular}{lcccccc}
\toprule
Model & Read ms & Prep ms & H2D ms & Forward ms & Write ms & E2E ms \\
\midrule
B0 fusion & 4.89 $\pm$ 0.64 & 2.80 $\pm$ 0.13 & 0.32 $\pm$ 0.02 & 27.43 $\pm$ 0.45 & 0.67 $\pm$ 0.09 & 36.81 $\pm$ 1.33 \\
B1 fusion & 4.49 $\pm$ 0.08 & 2.74 $\pm$ 0.05 & 0.31 $\pm$ 0.01 & 24.59 $\pm$ 0.61 & 0.52 $\pm$ 0.02 & 33.31 $\pm$ 0.71 \\
\bottomrule
\end{tabular}
\end{table}

\begin{table}[H]
\centering
\footnotesize
\caption{Fusion diagnostics for the dual-regression post-processing rule. The oracle-match rate measures how often the final fusion choice matches the lower-error branch for visible joints. Values are mean $\pm$ std over five seeds. The table includes all currently available diagnostic exports, including a shared-10 row for \texttt{B0}.}
\label{tab:fusion_diagnostics}
\begin{tabular}{llccccc}
\toprule
Model & Eval set & HM sel. & Coord sel. & $\alpha$ mean & Disagr. mean & Oracle match \\
\midrule
B0 & native-21 & 0.426 $\pm$ 0.025 & 0.574 $\pm$ 0.025 & 0.065 $\pm$ 0.007 & 0.163 $\pm$ 0.022 & 0.481 $\pm$ 0.012 \\
B0 & shared-10 & 0.407 $\pm$ 0.024 & 0.593 $\pm$ 0.024 & 0.065 $\pm$ 0.007 & 0.162 $\pm$ 0.018 & 0.474 $\pm$ 0.011 \\
B1 & native-10 & 0.501 $\pm$ 0.053 & 0.499 $\pm$ 0.053 & 0.084 $\pm$ 0.003 & 0.171 $\pm$ 0.018 & 0.529 $\pm$ 0.037 \\
\bottomrule
\end{tabular}
\end{table}

\begin{table}[H]
\centering
\footnotesize
\caption{Hyperparameter ablation relative to the baseline \texttt{B0} configuration, using native-21 fused accuracy on RHD and fused end-to-end benchmark latency. Values are mean $\pm$ std over five seeds.}
\label{tab:hyper_ablation}
\begin{tabular}{llcccc}
\toprule
Exp & Change vs. B0 & SSE $\downarrow$ & EPE $\downarrow$ & PCK@0.2 $\uparrow$ & E2E ms $\downarrow$ \\
\midrule
B0 & baseline & 0.3098 $\pm$ 0.0236 & 0.0779 $\pm$ 0.0054 & 0.9112 $\pm$ 0.0071 & 36.81 $\pm$ 1.33 \\
L1 & $\lambda_{hm}=1.5$, $\lambda_{coord}=0.5$ & 0.3583 $\pm$ 0.0534 & 0.0725 $\pm$ 0.0096 & 0.9081 $\pm$ 0.0092 & 35.77 $\pm$ 0.45 \\
L2 & $\lambda_{hm}=0.5$, $\lambda_{coord}=1.5$ & 0.3073 $\pm$ 0.0260 & 0.0812 $\pm$ 0.0022 & 0.9115 $\pm$ 0.0095 & 35.75 $\pm$ 0.49 \\
S1 & $\sigma=1.5$ & 0.3252 $\pm$ 0.0381 & 0.0707 $\pm$ 0.0055 & 0.9116 $\pm$ 0.0086 & 35.72 $\pm$ 0.25 \\
S2 & $\sigma=3.0$ & 0.3137 $\pm$ 0.0488 & 0.0833 $\pm$ 0.0058 & 0.9177 $\pm$ 0.0081 & 36.48 $\pm$ 0.53 \\
W1 & $w=5.0$ & 0.3213 $\pm$ 0.0384 & 0.0737 $\pm$ 0.0058 & 0.9150 $\pm$ 0.0099 & 36.21 $\pm$ 0.39 \\
W2 & $w=15.0$ & 0.2979 $\pm$ 0.0441 & 0.0750 $\pm$ 0.0031 & 0.9140 $\pm$ 0.0124 & 35.91 $\pm$ 0.43 \\
E1 & $\epsilon=1.0$ & 0.2994 $\pm$ 0.0591 & 0.0695 $\pm$ 0.0011 & 0.9175 $\pm$ 0.0113 & 35.84 $\pm$ 0.23 \\
E2 & $\epsilon=4.0$ & 0.3106 $\pm$ 0.0106 & 0.0873 $\pm$ 0.0056 & 0.9110 $\pm$ 0.0067 & 36.13 $\pm$ 0.52 \\
\bottomrule
\end{tabular}
\end{table}

\end{document}
